
\chapter{Hyper-Parameters and Training Settings}
\label{app:hyper}
The distillation hyper-parameters ($\alpha = \beta = 0.4$, $\gamma = 0.2$, $\delta = 0.3$, $T = 4$,
$\tau = 1$) are carried over from the Phase-2 report's Table 3.1 unchanged, so that the audit tests
the method as proposed; Section~\ref{sec:5.3} sweeps each of them. Students train for 6 epochs at $3
\times 10^{-5}$ ($10^{-3}$ for the BiLSTM), batch 32, 10 per cent warm-up, linear decay, gradient
clipping at 1.0, early stopping with patience 2 on validation macro-F1; the matched-steps controls
disable early stopping and run 13 and 35 epochs. Teachers train for 5 epochs at $2 \times 10^{-5}$,
batch 32. Maximum sequence length 128.

\begin{table}[htbp]
\centering
\small
\caption{Training settings of the teachers and the students.}
\label{tab:hyper}
\begin{tabular}{p{3.6cm}p{4.2cm}p{5.8cm}}
\toprule
Setting & Teachers & Students \\
\midrule
Epochs & 5 & 6; 13 and 35 for the matched-steps controls \\
Learning rate & $2 \times 10^{-5}$ & $3 \times 10^{-5}$; $10^{-3}$ for the BiLSTM \\
Batch size & 32 & 32 \\
Schedule & 10\% warm-up, linear decay & 10\% warm-up, linear decay \\
Gradient clipping & 1.0 & 1.0 \\
Model selection & best validation macro-F1 epoch & best validation macro-F1 epoch; early stopping with patience 2, off for the matched-steps controls \\
Maximum length & 128 tokens & 128 tokens \\
Precision & mixed, on GPU & mixed, on GPU \\
Distillation & -- & $\alpha = \beta = 0.4$, $\gamma = 0.2$, $\delta = 0.3$, $T = 4$, $\tau = 1$ \\
Seeds & one run per teacher & 3; 1 for ablations, sweeps and controls \\
\bottomrule
\end{tabular}
\end{table}
